\documentclass[11pt]{article}
\usepackage[margin=1in]{geometry}
\usepackage{amsmath,amssymb,amsthm,mathtools}
\usepackage{enumitem}
\usepackage{microtype}
\usepackage[hidelinks]{hyperref}
\numberwithin{equation}{section}
\newtheorem{theorem}{Theorem}[section]
\newtheorem{proposition}[theorem]{Proposition}
\newtheorem{lemma}[theorem]{Lemma}
\newtheorem{corollary}[theorem]{Corollary}
\theoremstyle{definition}
\newtheorem{definition}[theorem]{Definition}
\theoremstyle{remark}
\newtheorem{remark}[theorem]{Remark}
\newcommand{\E}{\mathbb E}
\newcommand{\Pp}{\mathbb P}
\newcommand{\Var}{\operatorname{Var}}
\newcommand{\TV}{\mathrm{TV}}
\newcommand{\Unif}{\mathrm{Unif}}
\newcommand{\dTV}{d_{\TV}}
\newcommand{\nuUnif}{\nu_{\Unif}}
\newcommand{\1}{\mathbf 1}
\newcommand{\Z}{\mathbb Z}
\newcommand{\cA}{\mathcal A}
\newcommand{\cI}{\mathcal I}
\newcommand{\cT}{\mathcal T}
\newcommand{\cU}{\mathcal U}
\newcommand{\Hh}{\mathsf H}
\DeclareMathOperator{\supp}{supp}
\DeclareMathOperator{\dist}{dist}
\DeclareMathOperator{\Cov}{Cov}
\title{Time-Inhomogeneous Riffle Shuffling}
\author{}
\date{}
\makeatletter
\renewcommand{\maketitle}{\begin{center}{\LARGE\@title\par}\end{center}\vspace{1em}}
\makeatother
\begin{document}
\maketitle
\begin{abstract}
We study riffle shuffles with deterministic pile sizes whose distribution may
vary on the time scale $\log N$. We prove uniform upper and lower bounds in
total variation in terms of two quantities: the size of the largest word
class and a variational expression describing cold spots. A fixed negative
margin gives mixing, whereas a fixed positive margin gives distance tending
to one, under the stated local nondegeneracy conditions. The estimates are
first obtained directly from the finite sequence of pile sizes. They then
transfer to time-dependent profiles under scalar integrability and local
empirical approximation. The upper bound also applies to random pile sizes
by averaging. The proofs use exact sampling with prescribed column counts,
conditional exploration of the shuffle graph, and an ascent test on
localized collision sets.
\end{abstract}

\section{Introduction}

A riffle shuffle cuts a deck into ordered piles and interleaves them
uniformly, preserving the order within each pile. Its inverse has a useful
description: assign each card a pile label, uniformly subject to the pile
sizes, and sort by that label while preserving ties. Iterating this
description encodes the shuffle by a matrix whose columns are independent
and have prescribed counts. Equal row words record the part of the initial
order which has survived the shuffles.

The distribution of the pile sizes need not be constant in time. A deck may
first undergo nearly balanced cuts, then strongly asymmetric cuts, or pass
through intervals with different empirical distributions. In this setting,
the time at which a group of cards is formed matters as well as its final
size. The analysis therefore involves costs accumulated over subintervals
of the shuffling history.

\subsection{Main results}

Fix $k\ge2$ and $N\ge2$. All logarithms are natural. A cut array
consists of nonnegative integers $n_i^{(j)}$ with
$\sum_{i=0}^{k-1}n_i^{(j)}=N$. Write
$p^{(j)}=(n_0^{(j)},\ldots,n_{k-1}^{(j)})/N$ for its pile proportions.
Let $P_N^K$ denote the law of the inverse of the chronological
composition of the first $K$ shuffles, started from the identity deck,
and let $\nuUnif$ be uniform on permutations of $N$ cards.
Inversion preserves total variation, so every bound below also holds
for the forward shuffle. We write $\dTV$ for total-variation distance.

For $p\in D_k=\{p\in[0,1]^k:\sum_i p_i=1\}$, write
\[
 a_p=\log\frac1{\max_i p_i},\qquad
 \psi_p(\theta)=-\log\sum_i p_i^\theta.
\]
A time profile $(\mu_t)$ is a family of probability measures on $D_k$.
Let $A(r,s)=\int_r^s\E_{\mu_t}a_p\,dt$ and
$B_\theta(r,s)=\int_r^s\E_{\mu_t}\psi_p(\theta)\,dt$.
The two relevant margins are
\begin{align*}
 \Lambda_0(T)&=1-A(0,T),\\
 \Lambda_1(T)&=\sup_{0\le s\le u\le T}\inf_{\theta\in[3,4]}
 \left\{\frac{\theta+3}{2}\bigl(1-A(0,s)\bigr)
       -B_\theta(s,u)-2B_2(u,T)\right\}.
\end{align*}

A column is $g$-good if two distinct piles each contain at least
$gN$ cards. The local $(g,\gamma,\rho)$-good condition on $[r,T]$
requires at least $\gamma\rho\log N$ good columns with
$j/\log N\in[a,a+\rho]$, for every real $a$ with
$[a,a+\rho]\subseteq[r,T]$.
Regular partitions, scalar integrability, and almost-likeness mean
\eqref{eq:partition-regularity}, \eqref{eq:scalar-integrability}, and
\eqref{eq:almost-like}, respectively. The same selected partition is
used in these conditions. Almost-likeness asks that its cell frequencies
on each $\rho$-grid interval approximate the integrated profile masses,
with error at most $\varphi$ times the interval length plus $1/\log N$.

Our main upper and lower bounds apply on opposite sides of the zero
margin. Their proofs are given in Sections~\ref{sec:upper-bound}
and~\ref{sec:lower-bound}.

\begin{theorem}[Buffered profile upper bound]\label{thm:upper-profile}
Fix $k\ge2$, $T>0$ and $\eta>0$. There exists $\chi_0>0$ such that,
for every $0<\chi\le\chi_0$ and every regular Borel partition $\mathcal P$ of
mesh $\chi$ with regularity constant $1$, the following holds.
For every $0<g<1/k$ and $\gamma>0$, there exists $\rho_0>0$ such
that each $0<\rho\le\rho_0$ admits $\varphi_0>0$. Fix
$0<\varphi\le\varphi_0$. Let $\mu$ satisfy the scalar-integrability
condition for this partition on $[0,T]$, and suppose
\[
 \max\{\Lambda_0(T),\Lambda_1(T)\}\le-\eta.
\]
For every $\varepsilon>0$, all sufficiently large $N$ and every array
with $K=\lfloor T\log N\rfloor$ which is
$(\mathcal P,\rho,\varphi)$-almost like $\mu$ on $[0,T]$ and locally
$(g,\gamma,\rho)$-good on $[0,T]$ satisfy
\[
 \dTV(P_N^K,\nuUnif)\le\varepsilon.
\]
The partition and all meshes are fixed before $\mu$ and $N$.
\end{theorem}

\begin{theorem}[Macroscopic lower bound]\label{thm:macroscopic-lower}
Fix $k\ge2$, $T>0$ and $\eta>0$. There are $0<\chi<1/k$, a $\chi$-regular finite Borel
partition $\mathcal R$ with regularity constant $1$, and $\rho_0>0$ such that
\[
 2\rho_0\log k\le\frac1{100}\min\{\eta/200000,1/1000\}.
\]
For every $0<\rho\le\rho_0$ there exists $\varphi_0>0$ with the following
property. Fix $0<\varphi\le\varphi_0$, $\chi_g,\gamma>0$, and a time
profile $\mu$ satisfying the scalar integrability condition~\eqref{eq:scalar-integrability} on $[0,T]$ for $\mathcal R$. If
\[
 \max\{\Lambda_0(T),\Lambda_1(T)\}\ge\eta,
\]
then, for every $\varepsilon>0$, there is $N_0\ge2$ such that all
$N\ge N_0$ and all cut arrays with $K=\lfloor T\log N\rfloor$ which are
$(\mathcal R,\rho,\varphi)$-almost like $\mu$ on $[0,T]$ and satisfy the local
good-column condition with parameters $\chi_g,\gamma,\rho$ on $[0,T]$
have $\dTV(P_N^K,\nuUnif)\ge1-\varepsilon$. In particular the distance
tends to one for every sequence of arrays eventually satisfying these
fixed hypotheses, even if the applicable empirical alternative varies
with $N$.
\end{theorem}

The finite-array versions of these results use the actual sequence of
pile sizes. Put $S_N=K/\log N$ and let $p_N(t)=p^{(j)}$ on
$((j-1)/\log N,j/\log N)$. Define
\[
 A^{(N)}(r,s)=\int_r^s a_{p_N(t)}\,dt,\qquad
 B^{(N)}_\theta(r,s)=\int_r^s\psi_{p_N(t)}(\theta)\,dt
 \quad(0\le r\le s\le S_N),
\]
and
\begin{align*}
 \beta^{(N)}(s)&=1-A^{(N)}(0,s),\\
 \Lambda_1^{(N)}&=\sup_{0\le s\le u\le S_N}\inf_{\theta\in[3,4]}
 \left\{\frac{\theta+3}{2}\beta^{(N)}(s)
       -B^{(N)}_\theta(s,u)-2B^{(N)}_2(u,S_N)\right\}.
\end{align*}
For the upper bound, empirical margin at least $\eta$ means
$\beta^{(N)}(S_N)\le-\eta$ and that the infimum in this display
is at most $-\eta$ for every grid pair
$s=i/\log N$, $u=j/\log N$, $0\le i\le j\le K$.
The lower bound uses the displayed supremum over all real endpoints.

\begin{theorem}[Empirical upper bound]\label{thm:upper-empirical}
Fix $k\ge2$ and $T,\eta_0,g_0,\gamma_0>0$. There exists $\rho_0>0$
such that, for every fixed $0<\rho\le\rho_0$, there exist $c>0$ and
$N_0\ge2$ with the following property. For $N\ge N_0$ and
$K\le T\log N$, suppose the deterministic array on $N$ labels has
empirical margin at least $\eta_0$ and satisfies the local
$(g_0,\gamma_0,\rho)$-good condition on $[0,K/\log N]$. Then
\begin{equation}\label{eq:upper-empirical-tv}
 \dTV(P_N^K,\nuUnif)\le N^{-c}.
\end{equation}
The choices of $c,N_0$ precede the cut array and the horizon $K$.
\end{theorem}

\begin{theorem}[Empirical lower bound]\label{thm:empirical-lower}
Fix $k\ge2$, $T>0$, positive margins $\eta_0,\eta_1$ and arbitrary real parameters
$\chi_g,\gamma,\rho$. For every $\varepsilon>0$, there is $N_0\ge2$ such
that all $N\ge N_0$ and all deterministic cut arrays with $K=\lfloor T\log N\rfloor$ have
$\dTV(P_N^K,\nuUnif)\ge1-\varepsilon$ whenever either
\[
 \beta^{(N)}(S_N)\ge\eta_0,
\]
or all of the following hold:
\[
 \chi_g,\gamma,\rho>0,\qquad
 2\rho\log k\le\frac1{100}\min\{\eta_1/100000,1/1000\},\qquad
 \Lambda_1^{(N)}\ge\eta_1,
\]
and the cut array satisfies the local good-column condition with
parameters $\chi_g,\gamma,\rho$ on $[0,T]$.
\end{theorem}

\subsection{Consequences and discussion}

The macroscopic upper and lower bounds retain a fixed buffer away from
zero. They do not require a measurable probability kernel in time. Instead,
they require integrability of the scalar functions which enter the
argument: the cell probabilities of the selected partition, the maximal
cost, and the power-sum costs for $\theta\in[2,4]$.

\begin{corollary}[Critical-time upper bound]\label{cor:upper-critical}
Suppose a consistent infinite sequence of cuts is specified for each
$N$, and these arrays admit arbitrarily fine approximations to $\mu$:
for each $T>0$ and each mesh bound $\chi_*>0$ there are
$0<\chi\le\chi_*$, a $\chi$-regular partition with constant $1$ satisfying
scalar integrability, and $g,\gamma,\bar\rho>0$, with $g<1/k$,
such that for every $0<\rho\le\bar\rho$ and every $\varphi>0$ the
arrays at horizon $\lfloor T\log N\rfloor$ are eventually almost
like $\mu$ and locally good with these parameters. For every $t\ge0$ with
\[
 t>\inf\{T>0:\max\{\Lambda_0(T),\Lambda_1(T)\}<0\},
 \qquad \inf\varnothing=+\infty,
\]
we have $\dTV(P_N^{\lfloor t\log N\rfloor},\nuUnif)\to0$.
Moreover, for every $0<\varepsilon<1$,
\[
 \limsup_{N\to\infty}\frac{t_{\rm mix}^{(N)}(\varepsilon)}{\log N}
       \le\inf\{T>0:\max\{\Lambda_0(T),\Lambda_1(T)\}<0\}.
\]
Here $t_{\rm mix}^{(N)}(\varepsilon)$ is the infimum of the integers
$m\ge0$ for which $\dTV(P_N^m,\nuUnif)\le\varepsilon$.
Both infima are interpreted in $[0,+\infty]$, with empty infimum
$+\infty$.
\end{corollary}

The lower bound applies at every time with a strictly positive margin
and the specified approximation conditions.

Uniformity in the cut array also gives the random-pile upper bound. After
conditioning on the pile sizes, the empirical estimate applies on the
admissible event. The total-variation distance of the mixture is at most
the probability of its complement plus the deterministic error. No
independence between the pile sizes at different times is needed for this
averaging step.

\subsection{Proof techniques and challenges}

The organization follows the first-moment, exponential-moment, and
cold-spot approach to riffle shuffling in~\cite{SSW}. The fixed-column model
requires care with dependence between cards. Conditional on the cards
already exposed, the remaining labels are sampled without replacement;
one cannot replace them by independent draws with the original pile
probabilities.

For the upper bound, we represent the inverse shuffle by sorting an
independent uniform permutation inside the components of its equality
graph. A truncated second-moment estimate reduces mixing to control of
common edges of two independent graphs. A good initial column supplies a
reserve of labels. Endpoint fluctuation estimates and a classification by
empirical digit costs then give a negative exponent for the first moment.
The same argument is applied inside the actual reservoirs left by an
exploration. Conditional moment estimates yield the exponential moment
needed by the graph comparison.

For the lower bound, successive hypergeometric thinning controls prefix
counts. The cold margin produces a finite tilted type class and a
deterministic union of short intervals. Collision estimates show that many
equality edges lie in a translate of this set. Sorting inside graph
components creates an excess of ascents there. A concentration estimate for
component sorting and a uniform bound over translates distinguish the
shuffle from a uniform permutation. The translate is selected from the
graph, before the independent uniform permutation used in the comparison.

We give the setup and finite probabilistic tools in Section~2, the upper
bound and its consequences in Section~3, the first-moment estimate in
Section~4, and the conditional exploration in Section~5. Section~6 proves
the lower bound.
